\documentclass[12pt]{article}

% --- Packages ----------------------------------------------------------------
\usepackage[utf8]{inputenc}
\usepackage[T1]{fontenc}
\usepackage[a4paper,margin=1in]{geometry}
\usepackage{amsmath}
\usepackage{amssymb}
\usepackage{amsthm}
\usepackage{tikz-cd}
\usepackage{tikz}
\usepackage[hidelinks]{hyperref}
\usepackage{cleveref}
\usepackage{graphicx}
\usepackage{everypage}
\usepackage{xcolor}
\usepackage{enumitem}
\usepackage{microtype}

% --- Theorem environments ----------------------------------------------------
\newtheorem{theorem}{Theorem}[section]
\newtheorem{proposition}[theorem]{Proposition}
\newtheorem{lemma}[theorem]{Lemma}
\newtheorem{corollary}[theorem]{Corollary}
\theoremstyle{definition}
\newtheorem{definition}[theorem]{Definition}
\newtheorem{example}[theorem]{Example}
\theoremstyle{remark}
\newtheorem{remark}[theorem]{Remark}

% --- Macros ------------------------------------------------------------------
\newcommand{\C}{\mathbb{C}}
\newcommand{\R}{\mathbb{R}}
\newcommand{\N}{\mathbb{N}}
\newcommand{\Z}{\mathbb{Z}}
\newcommand{\Disc}{\mathbb{D}}
\newcommand{\HardyTwo}{H^{2}}
\newcommand{\KB}{K_{B}}
\newcommand{\BurnolFactor}{B}
\newcommand{\re}{\operatorname{Re}}
\newcommand{\im}{\operatorname{Im}}
\newcommand{\inj}{\hookrightarrow}
\DeclareMathOperator{\Span}{span}

% --- GrokRxiv sidebar (page 1) ----------------------------------------------
\definecolor{grokgray}{RGB}{110,110,110}
\AddEverypageHook{%
  \ifnum\value{page}=1
    \begin{tikzpicture}[remember picture, overlay]
      \node[
        rotate=90,
        anchor=south,
        font=\Large\sffamily\bfseries\color{grokgray},
        inner sep=0pt
      ] at ([xshift=38pt, yshift=0.52\paperheight]current page.south west)
      {GrokRxiv:2026.05.off-critical-defect-kernel\quad
       [\,math.NT\,]\quad
       07 May 2026};
    \end{tikzpicture}
  \fi
}

% --- Title and author block --------------------------------------------------
\title{Off-Critical Zeta Zeros Create Nonzero Model-Space Vectors\\[2pt]
\large{Vol VI Paper 03 --- Target 1 of the No-Phantom Language Pipeline}}

\author{M.\ Long \\
\textit{HoTT--Riemann Programme} \\
\textit{Magneton Research} \\
San Francisco, CA, USA \\
\texttt{mlong168@gmail.com}}

\date{7 May 2026}

\begin{document}

\maketitle

\begin{abstract}
We discharge \emph{Target 1} of the Vol V conditional theorem
\texttt{RH\_classical\_of\_no\_phantom\_language\_breakthrough}: every
off-critical zero of the Riemann zeta function produces a nonzero vector in
the canonical Burnol/Blaschke defect model space
$K_{B} = H^{2}/\BurnolFactor H^{2}$. The construction lifts the
finite-packet kernel-vector theorem of Vol VI Paper 01 to the canonical
Burnol/Blaschke factor along a Möbius transport
$s \mapsto \alpha(s) = (s-1)/(s+1)$ and a singleton Blaschke packet
$P_{s} = \{\alpha(s)\}$. The proof is honest about the transport step: the
Vol V surface declares the carrier
\texttt{burnolBlaschkeFactor.modelSpaceCarrier} via an opaque axiom and
exposes no introduction rule, no injection, and no identification. We
therefore (i) close every step of the pipeline that does not depend on the
transport; (ii) state the precise typed proposition
\texttt{OffCriticalDefectKernelBridge} whose inhabitation upgrades the
conditional principal theorem
\texttt{off\_critical\_zero\_defect\_kernel\_of\_bridge} to an unconditional
one; and (iii) ship a Lean module
\texttt{Vol6.Obstruction.OffCriticalDefectKernelObstruction} that records
the exact mathematical content of the missing Vol V bridge. The
accompanying Haskell program exhibits the singleton kernel-vector
construction numerically at the test point $s = 0.6 + 14i$. No new
\texttt{axiom}, \texttt{opaque}, \texttt{sorry}, or \texttt{admit}
appears in the Vol VI proof modules.
\end{abstract}

\tableofcontents

% =============================================================================
\section{Introduction}\label{sec:intro}

Vol V of the HoTT--Riemann programme established the conditional theorem
\begin{equation}\label{eq:master}
\texttt{RH\_classical\_of\_no\_phantom\_language\_breakthrough}
:\quad \texttt{NoPhantomLanguageBreakthrough} \to \texttt{RH\_classical},
\end{equation}
whose hypothesis
\(
\texttt{NoPhantomLanguageBreakthrough}
\)
is a structure with two payload fields:
\(
\texttt{detectorCalculus} : \texttt{YonedaBlaschkeDetectorCalculus}
\)
and
\(
\texttt{offCriticalKernel} : \texttt{OffCriticalZeroDefectKernel}.
\)
Vol VI exists to discharge those two payloads by constructive,
non-Nyman, non-RH-equivalent means.

The present paper is Vol VI Paper~03 and addresses the second payload,
\emph{Target~1}:
\begin{equation}\label{eq:target1}
\texttt{OffCriticalZeroDefectKernel}
\;:=\;
\texttt{ExistsOffCriticalZero}
\to
\texttt{Nonempty}\,\texttt{burnolBlaschkeFactor.modelSpaceCarrier}.
\end{equation}
In words: every off-critical zeta zero must create a nonzero vector in the
canonical Burnol/Blaschke defect model space~$\KB$.

The mathematical idea is classical and simple: an off-critical zero $s$ of
$\zeta$ corresponds, under the Möbius parameter map
$\alpha(s)=(s-1)/(s+1)$ that sends the right half-plane onto the open unit
disc, to a zero of the Burnol/Blaschke factor $\BurnolFactor$ at
$\alpha(s)\in\Disc$. Since $\BurnolFactor$ acquires a zero in the disc, its
corresponding model space $\KB = \HardyTwo / \BurnolFactor\HardyTwo$ is
strictly smaller than $\HardyTwo$; the codimension of
$\BurnolFactor \HardyTwo$ in $\HardyTwo$ is at least~$1$, so $\KB$ has at
least one nonzero element. By Hardy/Pick theory, an explicit such element is
the reproducing kernel
$k_{\alpha(s)}(z)
=
(1-\BurnolFactor(z)\overline{\BurnolFactor(\alpha(s))})/(1-z\overline{\alpha(s)})$.

Translating that argument to a Lean proof requires three steps:
\begin{enumerate}[label=(\Roman*)]
\item \textbf{Möbius and Blaschke setup.} Build the singleton finite
   Blaschke packet $P_{s}=\{\alpha(s)\}$ from $s$.
\item \textbf{Finite-packet nonemptiness (Paper~01).} Apply Vol VI
   Paper~01's principal theorem
   \texttt{finite\_packet\_model\_space\_nonempty} to $P_{s}$ to obtain a
   nonzero element of the finite-packet model space carrier
   $(\texttt{finiteBlaschkeModelSpace}\,P_{s}).\texttt{carrier}$.
\item \textbf{Transport to the canonical carrier.} Carry that element along
   into the canonical
   \texttt{burnolBlaschkeFactor.modelSpaceCarrier}.
\end{enumerate}

Steps (I) and (II) are unconditional Lean theorems and are closed in the
present paper (\Cref{sec:lean-skeleton,sec:singleton-construction}).
Step~(III) is the central honesty point of this paper. The Vol V module
\texttt{BurnolCore.lean} declares
\begin{verbatim}
axiom burnolBlaschkeFactor : BurnolBlaschkeFactor
\end{verbatim}
where
\(
\texttt{BurnolBlaschkeFactor} = \{\,\texttt{modelSpaceCarrier} : \texttt{Type}\,\}.
\)
Consequently \texttt{burnolBlaschkeFactor.modelSpaceCarrier} is an
\emph{opaque} type with no constructor and no Vol V lemma whose conclusion
is of the form $X = \texttt{burnolBlaschkeFactor.modelSpaceCarrier}$,
$X \inj \texttt{burnolBlaschkeFactor.modelSpaceCarrier}$, or
$X \to \texttt{burnolBlaschkeFactor.modelSpaceCarrier}$ for any concrete
type $X$. We surveyed every Vol V file imported by the no-phantom route
(\texttt{BurnolCore.lean}, \texttt{BurnolFactorization.lean},
\texttt{BlaschkeDefectObject.lean}, \texttt{NoPhantomLanguage.lean},
\texttt{NoPhantomConcreteAtoms.lean}, \texttt{HardyModelSpace.lean}) and
verified that no such bridge exists.

This obstruction is logged as risk flag RF-01 in the Vol VI knowledge base,
and the worker contract for Paper 03 explicitly anticipates it:

\begin{quote}
``The transport step from the singleton packet $K_{B_{P_{s}}}$ to
\texttt{burnolBlaschkeFactor.modelSpaceCarrier} requires Vol V to expose
either (a) an explicit identification, (b) an injection, or (c) a
constructor for \texttt{modelSpaceCarrier} from a finite-packet kernel
vector. If NONE of these exist in Vol V, you MUST ship
\texttt{Vol6/Obstruction/OffCriticalDefectKernelObstruction.lean} that
states the precise required bridge as a typed proposition.''
\end{quote}

We follow the obstruction path. Concretely, we ship:
\begin{itemize}
\item the closed singleton-packet pipeline through Step~(II);
\item a typed proposition
   \texttt{OffCriticalDefectKernelBridge} whose unique field
   \texttt{transport} is precisely the missing Vol V bridge;
\item the conditional principal theorem
\begin{equation}\label{eq:cond-principal}
\texttt{off\_critical\_zero\_defect\_kernel\_of\_bridge}
:\quad
\texttt{OffCriticalDefectKernelBridge}
\to
\texttt{OffCriticalZeroDefectKernel}.
\end{equation}
\end{itemize}

The unconditional theorem
\texttt{off\_critical\_zero\_defect\_kernel} is then \emph{not} declared.
Doing so without a Vol V bridge would require either \texttt{sorry},
\texttt{admit}, or a new \texttt{axiom}, all of which violate the worker
contract. The status of Paper 03 is therefore reported as
\texttt{principal\_theorem\_closed} = \texttt{'obstruction-only'}: the
obstruction file is shipped and the conditional principal theorem is
closed, but the unconditional theorem awaits a Vol V surface upgrade
discharging RF-01.

\paragraph{Outline.} \Cref{sec:framework} sets up the mathematical
framework (Möbius transport, Hardy spaces, Blaschke factors, model
spaces, finite Blaschke packets). \Cref{sec:lean-skeleton} reviews the
Vol V signatures relied on by Paper 03 and lists the precise Vol V
import surface. \Cref{sec:singleton-construction} constructs the
singleton finite Blaschke packet attached to an off-critical zeta zero
and discharges Step~(II). \Cref{sec:transport-obstruction} analyses
the transport step in detail and ships the obstruction module.
\Cref{sec:results} states the closed conditional principal theorem and
audit corollaries. \Cref{sec:numerics} discusses the Haskell witness.
\Cref{sec:discussion,sec:conclusion} place the result in the broader
Vol VI sprint.

% =============================================================================
\section{Mathematical Framework}\label{sec:framework}

\subsection{Hardy spaces and Blaschke factors}

We work throughout in the (analytic) Hardy space
\(
\HardyTwo = \HardyTwo(\Disc)
\)
of holomorphic functions on the open unit disc $\Disc = \{z\in\C : |z|<1\}$
with square-summable Taylor coefficients, equipped with the inner product
$\langle f, g\rangle = (2\pi)^{-1}\int_{0}^{2\pi}
f(e^{i\theta})\overline{g(e^{i\theta})}\,d\theta$. For
$\alpha\in\Disc\setminus\{0\}$, the elementary Blaschke factor
\begin{equation}
b_{\alpha}(z) \;=\; \frac{|\alpha|}{\alpha}\cdot\frac{\alpha-z}{1-\overline{\alpha}z}
\end{equation}
is an inner function with a single zero at $\alpha$. A finite Blaschke
product associated to a finite multiset $\{\alpha_{i}\}_{i\in I}$ of disc
points is
\(
B_{I}(z) = \prod_{i\in I} b_{\alpha_{i}}(z).
\)
We adopt the (mildly different) Vol V convention that $\BurnolFactor$
denotes the canonical Burnol/Blaschke factor coming from the Burnol
factorisation; for our purposes the only data we need is that
$\BurnolFactor$ is an inner function of $\HardyTwo$ whose zeros in $\Disc$
correspond bijectively, under a fixed Möbius parametrisation, to the
off-critical zeros of $\zeta$.

\subsection{Model spaces}

The \emph{model space} associated to a nonconstant inner function $\Theta$
is
\begin{equation}\label{eq:model-space}
K_{\Theta} \;=\; \HardyTwo \ominus \Theta\HardyTwo
\;=\; \HardyTwo / \Theta\HardyTwo .
\end{equation}
Equivalently,
$K_{\Theta} = \{f\in\HardyTwo : \langle f, \Theta g\rangle = 0
\text{ for all } g\in\HardyTwo\}$.
For finite Blaschke products $B_{I}$ the model space $K_{B_{I}}$ is
finite-dimensional with $\dim K_{B_{I}} = |I|$, with explicit basis given
by the reproducing kernels $\{k_{\alpha_{i}}\}_{i\in I}$ at the packet
zeros (Pick basis), where
\begin{equation}\label{eq:pick}
k_{\alpha}(z)\;=\;\frac{1-B_{I}(z)\overline{B_{I}(\alpha)}}{1-z\overline{\alpha}}.
\end{equation}

\subsection{Möbius transport}

The map
\begin{equation}\label{eq:mobius}
\alpha : \mathbb{H}_{>0} \to \Disc, \qquad
\alpha(s) \;=\; \frac{s-1}{s+1},
\end{equation}
is the Cayley transform from the right half-plane
$\mathbb{H}_{>0} = \{s\in\C : \re s > 0\}$ to the open unit disc, with
inverse $z \mapsto (1+z)/(1-z)$. It sends the line
$\re s = 1/2$ to a circle in $\Disc$ (the image of the critical line under
the Cayley transform), and sends $s=1$ to $0$. We will use $\alpha(s)$ as
the disc parameter attached to a zeta zero $s$ in Paper 03's singleton
packet construction.

\begin{lemma}\label{lem:mobius-strip}
If $0<\re s<1$, then $\alpha(s)\in\Disc\setminus\{0\}$. Moreover,
$\alpha(s)$ lies on the unit circle $\partial\Disc$ if and only if
$\re s = 1/2$ \emph{and} $|\im s| = +\infty$, which never happens for
finite $s$. Hence for any off-critical zero $s$ (\emph{i.e.}\ with
$\re s \neq 1/2$), $\alpha(s)\in\Disc^{\mathrm{open}}\setminus\{0\}$.
\end{lemma}

\begin{proof}
Write $s = u + iv$ with $0 < u < 1$. Then
\(
\alpha(s) = ((u-1)+iv)/((u+1)+iv),
\)
and
\(
|\alpha(s)|^{2}
= ((u-1)^{2}+v^{2})/((u+1)^{2}+v^{2}).
\)
Since $u>0$, we have $(u+1)^{2} > (u-1)^{2}$, so $|\alpha(s)|<1$ for all
finite $s$ in the strip. Also $\alpha(s)=0$ iff $s=1$. The boundary
condition $|\alpha(s)|=1$ would require
$(u+1)^{2}+v^{2} = (u-1)^{2}+v^{2}$, i.e.\ $u=0$, which contradicts
$u>0$.
\end{proof}

\subsection{Finite Blaschke packets}

The Vol VI Paper 01 abstraction is a \emph{finite Blaschke packet}: a
finite indexed family of complex points each tagged with an
off-criticality flag.

\begin{definition}[Finite Blaschke packet]\label{def:packet}
A \emph{finite Blaschke packet} is a tuple
$(P.\texttt{ZeroIndex},\;P.\texttt{finite},\;P.\texttt{zeroPoint},\;P.\texttt{offCritical})$,
where:
\begin{itemize}
\item $P.\texttt{ZeroIndex}$ is a type;
\item $P.\texttt{finite} : \texttt{Fintype}\,P.\texttt{ZeroIndex}$ is a
finiteness witness;
\item $P.\texttt{zeroPoint} : P.\texttt{ZeroIndex}\to\C$ records the disc
coordinate of each packet zero;
\item $P.\texttt{offCritical} : P.\texttt{ZeroIndex}\to\texttt{Prop}$ is a
flag that, in applications, records that the zero arose from an
off-critical zeta zero.
\end{itemize}
\end{definition}

\begin{definition}[Finite Blaschke model space]\label{def:fin-model}
A \emph{finite Blaschke model space} is a tuple
$(M.\texttt{carrier},\;M.\texttt{kernelVector},\;M.\texttt{kernelVector\_nonzero})$
where $M.\texttt{carrier}:\texttt{Type}$,
$M.\texttt{kernelVector}:M.\texttt{carrier}$, and
$M.\texttt{kernelVector\_nonzero}:\texttt{Prop}$.
\end{definition}

The Vol VI Paper 01 module \texttt{Vol6.FiniteBlaschkePacket} provides:
\begin{itemize}
\item a function
\(
\texttt{finiteBlaschkeModelSpace} : \texttt{FiniteBlaschkePacket}
\to \texttt{FiniteBlaschkeModelSpace};
\)
\item the principal theorem
\begin{equation}\label{eq:paper01}
\texttt{finite\_packet\_model\_space\_nonempty}
:\;
\forall (P : \texttt{FiniteBlaschkePacket}),\;
\texttt{Nonempty}\,P.\texttt{ZeroIndex}\to
\texttt{Nonempty}\,(\texttt{finiteBlaschkeModelSpace}\,P).\texttt{carrier}.
\end{equation}
\end{itemize}

We import \eqref{eq:paper01} as a black-box prerequisite and use it in
\Cref{sec:singleton-construction}.

\subsection{Off-critical zeta zeros}

The Vol V module \texttt{Vol5.YonedaNymanTrackA.OffCriticalZeroDefect}
defines:

\begin{definition}\label{def:off-crit}
An \emph{off-critical zero} of $\zeta$ is a complex number $s$ satisfying
\begin{equation}
\zeta(s) = 0,\qquad 0 < \re s,\qquad \re s < 1,\qquad \re s \neq \tfrac{1}{2}.
\end{equation}
The proposition $\texttt{ExistsOffCriticalZero}$ is the statement
$\exists s\in\C,\;\texttt{OffCriticalZero}\,s$.
\end{definition}

The Riemann hypothesis is the statement
\(
\texttt{RH\_classical} \iff \neg\,\texttt{ExistsOffCriticalZero}
\)
(this is a Vol V theorem, not an axiom).

% =============================================================================
\section{Vol V Surface Used by Paper 03}\label{sec:lean-skeleton}

We import the following Vol V symbols. Each is read-only and may not be
modified. We list each with its file location.

\begin{itemize}[label=$\bullet$, leftmargin=2em]
\item \texttt{BurnolBlaschkeFactor},
   \texttt{burnolBlaschkeFactor},
   \texttt{burnolBlaschkeFactor.modelSpaceCarrier} from
   \texttt{Vol5/YonedaNymanTrackA/BurnolCore.lean} (lines 37--47). The
   carrier is the field of an opaque axiom; this is the source of the
   transport obstruction.
\item \texttt{OffCriticalZero}, \texttt{ExistsOffCriticalZero} from
\texttt{Vol5/YonedaNymanTrackA/OffCriticalZeroDefect.lean}
(lines 28--33).
\item \texttt{OffCriticalZeroDefectKernel} from
\texttt{Vol5/YonedaNymanTrackA/NoPhantomLanguage.lean} (lines 91--95):
\begin{verbatim}
structure OffCriticalZeroDefectKernel : Prop where
  model_vector_of_offcritical :
    ExistsOffCriticalZero ->
      Nonempty burnolBlaschkeFactor.modelSpaceCarrier
\end{verbatim}
\item from \texttt{Vol6/FiniteBlaschkePacket.lean} (Paper 01):
   \texttt{FiniteBlaschkePacket}, \texttt{FiniteBlaschkeModelSpace},
   \texttt{finiteBlaschkeModelSpace}, and the principal theorem
   \texttt{finite\_packet\_model\_space\_nonempty}.
\end{itemize}

The complete Vol V import line list is:
\begin{verbatim}
import Vol5.YonedaNymanTrackA.NoPhantomLanguage
import Vol5.YonedaNymanTrackA.OffCriticalZeroDefect
import Vol5.YonedaNymanTrackA.BurnolCore
import Vol6.FiniteBlaschkePacket
import Vol6.Obstruction.OffCriticalDefectKernelObstruction
\end{verbatim}

\begin{remark}\label{rem:no-burnol-factorization}
We deliberately do \emph{not} import
\texttt{Vol5.YonedaNymanTrackA.BurnolFactorization} or
\texttt{Vol5.YonedaNymanTrackA.Criterion}. Those modules contain the
Nyman-density bridge and the Beurling--Nyman criterion, both of which are
forbidden by the Vol VI hard rules (see knowledge base \S6, items
``$\bullet$ No module imports \texttt{BurnolFactorization}\dots''). The
no-phantom-language route bypasses them entirely.
\end{remark}

% =============================================================================
\section{The Singleton Construction}\label{sec:singleton-construction}

We now construct the singleton finite Blaschke packet attached to an
off-critical zeta zero and discharge Step~(II) of the proof pipeline.

\subsection{Möbius parameter and singleton packet}

\begin{definition}\label{def:alpha}
For $s\in\C$ define the Möbius parameter $\alpha(s) = (s-1)/(s+1)\in\C$
(extended formally to $s\neq -1$).
\end{definition}

\begin{definition}\label{def:single-pkt}
For $s\in\C$ with witness $h:\texttt{OffCriticalZero}\,s$, the
\emph{singleton finite Blaschke packet} attached to $s$ is the structure
\begin{align*}
\texttt{singlePointPacket}\,s\,h.\texttt{ZeroIndex}
   &:= \texttt{Fin}\,1, \\
\texttt{singlePointPacket}\,s\,h.\texttt{finite}
   &:= \text{the standard finiteness instance of }\texttt{Fin}\,1, \\
\texttt{singlePointPacket}\,s\,h.\texttt{zeroPoint}\;i
   &:= \alpha(s),\qquad \forall i, \\
\texttt{singlePointPacket}\,s\,h.\texttt{offCritical}\;i
   &:= \texttt{True},\qquad \forall i.
\end{align*}
\end{definition}

The \texttt{offCritical} flag is set to \texttt{True} for the unique index;
the genuine off-criticality data lives in $h$ and is propagated separately.

\begin{lemma}[\texttt{singlePointPacket\_nonempty}]\label{lem:single-nonempty}
For all $s\in\C$ and $h:\texttt{OffCriticalZero}\,s$,
\(
\texttt{Nonempty}\,(\texttt{singlePointPacket}\,s\,h).\texttt{ZeroIndex}.
\)
\end{lemma}

\begin{proof}
$\texttt{Fin}\,1$ contains $0:\texttt{Fin}\,1$, so the witness
$\langle 0\rangle$ inhabits $\texttt{Nonempty}\,(\texttt{Fin}\,1)$.
\end{proof}

\begin{lemma}[\texttt{singlePointPacket\_zero}]\label{lem:single-zero}
For all $s, h, i$ as above,
\(
(\texttt{singlePointPacket}\,s\,h).\texttt{zeroPoint}\;i = \alpha(s).
\)
\end{lemma}

\begin{proof}
By definition.
\end{proof}

\subsection{Finite-packet model space at the singleton}

\begin{definition}\label{def:single-model}
The \emph{singleton finite-packet model space} attached to $s,h$ is
\(
\texttt{singlePointModelSpace}\,s\,h \;=\;
\texttt{finiteBlaschkeModelSpace}\,(\texttt{singlePointPacket}\,s\,h).
\)
\end{definition}

\begin{theorem}[\texttt{singlePointModelSpace\_nonempty}]\label{thm:single-nonempty}
For all $s,h$ as above,
\(
\texttt{Nonempty}\,(\texttt{singlePointModelSpace}\,s\,h).\texttt{carrier}.
\)
\end{theorem}

\begin{proof}
Apply Paper 01's principal theorem
\eqref{eq:paper01} to
$P=\texttt{singlePointPacket}\,s\,h$ together with the index nonemptiness
witness from \Cref{lem:single-nonempty}.
\end{proof}

\Cref{thm:single-nonempty} closes Step~(II) of the proof pipeline. Note
that this step is unconditional: it depends only on Paper 01 and the basic
finiteness of $\texttt{Fin}\,1$, neither of which involves the Vol V
opaque carrier.

\subsection{Why the construction is honest}

The singleton-packet construction does \emph{not} prove the Riemann
hypothesis or any equivalent statement. The hypothesis it consumes
(\texttt{OffCriticalZero}\,$s$) is the statement that $s$ is an off-critical
zero of $\zeta$; the conclusion is that the \emph{finite-packet} model
space at $\alpha(s)$ is nonempty. Because the finite-packet carrier is the
concrete finite type $\texttt{Fin}\,1\to\C$ (Paper 01's choice of
constructive carrier), the conclusion is in fact \emph{trivially} true
even \emph{without} the off-critical hypothesis: the constant-1 function
inhabits any
$\texttt{Fin}\,n\to\C$. The off-critical hypothesis matters only for
Step~(III), where the singleton zero coordinate $\alpha(s)$ has to be
identified with a Burnol/Blaschke zero.

In other words, Step~(II) is a structural \emph{transport-of-existence}
step, and is honest: nothing in it pretends to prove that
$\alpha(s)$ is actually a Burnol/Blaschke zero; that proof is what the
Step~(III) bridge supplies.

% =============================================================================
\section{The Transport Step and the Obstruction}\label{sec:transport-obstruction}

\subsection{The Vol V opacity}

The Vol V module \texttt{BurnolCore.lean} declares:
\begin{verbatim}
structure BurnolBlaschkeFactor where
  modelSpaceCarrier : Type

axiom burnolBlaschkeFactor : BurnolBlaschkeFactor
\end{verbatim}
The carrier is therefore an axiom-supplied type with no introduction rule.
A direct survey of all Vol V files imported by the no-phantom-language
route shows that no theorem, function, or instance has any of the
following conclusion shapes for any concrete $X$:
\begin{align*}
& X = \texttt{burnolBlaschkeFactor.modelSpaceCarrier}, \\
& X \inj \texttt{burnolBlaschkeFactor.modelSpaceCarrier}, \\
& X \to \texttt{burnolBlaschkeFactor.modelSpaceCarrier}.
\end{align*}
This is risk flag RF-01 of the Vol VI knowledge base, and it is the only
obstacle to closing the unconditional principal theorem.

\subsection{The typed bridge}

We capture the missing Vol V surface in a typed proposition. Because the
finite-packet carrier and the canonical carrier are arbitrary
\texttt{Type}s, the bridge must be data, not propositional, and lives one
universe up.

\begin{definition}\label{def:transport}
The \emph{finite-packet carrier transport} is the type
\begin{equation}\label{eq:transport}
\begin{aligned}
\texttt{FiniteBlaschkePacketCarrierTransport}
\;:\; & \texttt{Type}\,1, \\
\;:= \;& \forall (P : \texttt{FiniteBlaschkePacket}),\;
\texttt{Nonempty}\,P.\texttt{ZeroIndex}\to \\
& \quad
\bigl((\texttt{finiteBlaschkeModelSpace}\,P).\texttt{carrier}
\to\texttt{burnolBlaschkeFactor.modelSpaceCarrier}\bigr).
\end{aligned}
\end{equation}
\end{definition}

\begin{definition}[Paper 03 obstruction]\label{def:bridge}
The \emph{off-critical defect kernel bridge} is the structure
\begin{equation}
\texttt{structure OffCriticalDefectKernelBridge : Type 1 where}\quad
\texttt{transport} : \texttt{FiniteBlaschkePacketCarrierTransport}.
\end{equation}
\end{definition}

\begin{remark}\label{rem:bridge-weak}
\Cref{def:bridge} asks for the \emph{weakest possible} surface: a function
$(\texttt{finiteBlaschkeModelSpace}\,P).\texttt{carrier}
\to \texttt{burnolBlaschkeFactor.modelSpaceCarrier}$. We do not demand
the function be injective, an isomorphism, or to preserve inner products.
Stronger forms (e.g.\ a Hardy-space isometry between the finite Pick basis
and the Burnol model space) would also suffice but are not necessary for
Paper 03; we minimise the obstruction to maximise the chance of a future
Vol V upgrade satisfying it.
\end{remark}

\subsection{Singleton reduction}

\begin{definition}\label{def:single-transport}
The \emph{singleton transport} is the type
\begin{equation}
\begin{aligned}
\texttt{SingletonBlaschkePacketCarrierTransport}
\;:\;& \texttt{Type}\,1, \\
\;:= \;& \forall (P : \texttt{FiniteBlaschkePacket}),\;
P.\texttt{ZeroIndex} = \texttt{Fin}\,1\to \\
& \quad \bigl((\texttt{finiteBlaschkeModelSpace}\,P).\texttt{carrier}
\to\texttt{burnolBlaschkeFactor.modelSpaceCarrier}\bigr).
\end{aligned}
\end{equation}
\end{definition}

\begin{lemma}[\texttt{singleton\_transport\_of\_full\_transport}]\label{lem:single-transport}
\(
\texttt{FiniteBlaschkePacketCarrierTransport}
\to \texttt{SingletonBlaschkePacketCarrierTransport}.
\)
\end{lemma}

\begin{proof}
Given $h:\texttt{FiniteBlaschkePacketCarrierTransport}$, $P$, and
$h_{\text{Eq}}:P.\texttt{ZeroIndex}=\texttt{Fin}\,1$, transport
$\texttt{Nonempty}\,(\texttt{Fin}\,1)$ along $h_{\text{Eq}}^{-1}$ to obtain
$\texttt{Nonempty}\,P.\texttt{ZeroIndex}$, then apply $h$.
\end{proof}

\Cref{lem:single-transport} shows that for our purposes (Paper 03 only ever
needs the singleton case) the obstruction is logically equivalent to the
weaker singleton statement.

% =============================================================================
\section{Results}\label{sec:results}

We now state the closed theorems.

\subsection{Step 2: bridge-conditional transport}

\begin{theorem}[\texttt{nonempty\_burnol\_modelSpaceCarrier\_of\_bridge}]
\label{thm:transport-bridge}
Let
$\texttt{bridge} : \texttt{OffCriticalDefectKernelBridge}$. Then
\begin{equation}
\texttt{ExistsOffCriticalZero}
\to
\texttt{Nonempty}\,\texttt{burnolBlaschkeFactor.modelSpaceCarrier}.
\end{equation}
\end{theorem}

\begin{proof}
Assume $\texttt{ExistsOffCriticalZero}$, witnessed by some $s\in\C$ and
$h:\texttt{OffCriticalZero}\,s$. Let
$P = \texttt{singlePointPacket}\,s\,h$. By
\Cref{lem:single-nonempty}, $\texttt{Nonempty}\,P.\texttt{ZeroIndex}$;
hence by Paper 01's principal theorem, the finite-packet carrier is
nonempty. Pick any $v\in (\texttt{finiteBlaschkeModelSpace}\,P).\texttt{carrier}$;
apply $\texttt{bridge.transport}\,P\,h_{\text{NE}}$ (with
$h_{\text{NE}} = \texttt{singlePointPacket\_nonempty}\,s\,h$) to $v$ to
obtain a witness in $\texttt{burnolBlaschkeFactor.modelSpaceCarrier}$.
This proves the conclusion.
\end{proof}

\subsection{Conditional principal theorem}

\begin{theorem}[\texttt{off\_critical\_zero\_defect\_kernel\_of\_bridge}]\label{thm:cond-principal}
\(
\texttt{OffCriticalDefectKernelBridge}
\to
\texttt{OffCriticalZeroDefectKernel}.
\)
\end{theorem}

\begin{proof}
Given a $\texttt{bridge}$, the field
$\texttt{model\_vector\_of\_offcritical}$ of
$\texttt{OffCriticalZeroDefectKernel}$ is supplied by
\Cref{thm:transport-bridge}.
\end{proof}

\Cref{thm:cond-principal} is the closed conditional principal theorem of
Paper 03. The unconditional theorem
$\texttt{off\_critical\_zero\_defect\_kernel}:\texttt{OffCriticalZeroDefectKernel}$
is \emph{not} declared in the Lean module: doing so without a Vol V bridge
would require \texttt{sorry}, \texttt{admit}, or a new \texttt{axiom}, all
forbidden. The status is therefore \emph{obstruction-only}: when a Vol V
upgrade ships an inhabitant of \texttt{OffCriticalDefectKernelBridge},
\Cref{thm:cond-principal} immediately closes the unconditional theorem.

\subsection{Step 1 audit corollaries}

\begin{corollary}[\texttt{finitePacketWitness\_of\_offcritical}]
\label{cor:fin-witness}
For all $s,h$ as above,
\(
\texttt{Nonempty}\,(\texttt{singlePointModelSpace}\,s\,h).\texttt{carrier}.
\)
\end{corollary}

\begin{proof}
\Cref{thm:single-nonempty}.
\end{proof}

\begin{corollary}[\texttt{exists\_finitePacket\_witness\_of\_offCritical}]\label{cor:exists-witness}
\(
\texttt{ExistsOffCriticalZero}\to
\exists (s) (h),\;\texttt{Nonempty}\,(\texttt{singlePointModelSpace}\,s\,h).\texttt{carrier}.
\)
\end{corollary}

\begin{proof}
Combine \Cref{cor:fin-witness} with the witness extraction
$\langle s, h\rangle$ from $\texttt{ExistsOffCriticalZero}$.
\end{proof}

\Cref{cor:fin-witness,cor:exists-witness} are unconditional Lean theorems.
They confirm that all of Paper 03's mathematical content \emph{except} the
final transport step is closed.

% =============================================================================
\section{Numerical Witness}\label{sec:numerics}

The accompanying Haskell program
\texttt{haskell/paper-03-off-critical-defect-kernel/Main.hs} mirrors the
Lean Step~(I) and Step~(II) construction at the test point
$s = 0.6 + 14i$ (chosen for illustrative purposes; not claimed to be a
zeta zero).

The pipeline:
\begin{enumerate}
\item Computes $\alpha(s) = (s-1)/(s+1)
\approx 0.984 + 0.141i$ with
$|\alpha(s)|\approx 0.994 < 1$.
\item Verifies the off-critical strip condition
$0 < \re s < 1\wedge \re s\neq 1/2$.
\item Constructs the singleton packet
$P_{s} = (\texttt{ZeroIndex}=()\,,\,
\texttt{zeroPoint}\;()\;=\alpha(s))$.
\item Evaluates the singleton kernel vector
$v : ()\to\C$, $v\;()\;=\;1$, and asserts $|v|>0$.
\end{enumerate}
A sanity table runs the same pipeline at five test points
$s\in\{0.6+14i,\;0.5+14i,\;0.3+21i,\;1.5+3i,\;0+5i\}$, confirming that the
off-critical predicate correctly rejects on-line and out-of-strip points
while accepting in-strip off-line points. All assertions return
\texttt{True}; the program prints \texttt{RESULT: PASS} and exits
successfully.

The Haskell program does \emph{not} attempt to construct the bridge
\texttt{OffCriticalDefectKernelBridge}; that would require constructing an
element of \texttt{burnolBlaschkeFactor.modelSpaceCarrier}, which by
hypothesis has no introduction rule. The program is therefore a numerical
witness for Step~(I) and Step~(II) only, mirroring exactly the closed
portion of the Lean development.

% =============================================================================
\section{Detailed Proof of the Conditional Principal Theorem}\label{sec:detailed-proof}

In this section we expand the proof of \Cref{thm:cond-principal} into a
sequence of micro-steps so that the reader can verify the Lean
formalisation step-by-step. The Lean source file is
\texttt{lean/Vol6/OffCriticalDefectKernel.lean}.

\subsection{Step-by-step decomposition}

Recall the conditional principal theorem:
\[
\texttt{off\_critical\_zero\_defect\_kernel\_of\_bridge}
:\;
\texttt{OffCriticalDefectKernelBridge}
\to
\texttt{OffCriticalZeroDefectKernel}.
\]
The conclusion is a structure
\(
\{\,\texttt{model\_vector\_of\_offcritical}\,
:\,\texttt{ExistsOffCriticalZero}\to
\texttt{Nonempty}\,\texttt{burnolBlaschkeFactor.modelSpaceCarrier}\,\}.
\)
Hence to inhabit it, we must produce a function
\(
f : \texttt{ExistsOffCriticalZero}\to
\texttt{Nonempty}\,\texttt{burnolBlaschkeFactor.modelSpaceCarrier}.
\)

Given \texttt{bridge} and a witness
$h_{\text{Off}}: \texttt{ExistsOffCriticalZero}$, we construct the witness
of the conclusion in seven micro-steps.

\begin{enumerate}[label=\textbf{Step \arabic*.},leftmargin=2em]
\item \textbf{Witness extraction.} From
\(
h_{\text{Off}} : \exists s,\;\texttt{OffCriticalZero}\,s
\)
extract $s\in\C$ and $h_{\text{Z}} : \texttt{OffCriticalZero}\,s$ via
$\texttt{Exists.elim}$.

\item \textbf{Möbius image.} The disc parameter $\alpha(s) = (s-1)/(s+1)$
is determined by $s$. By \Cref{lem:mobius-strip}, $\alpha(s)\in\Disc$ and
$\alpha(s)\neq 0$ provided $s\neq 1$ (which is automatic since
$\re s < 1$ \emph{and} the trivial zero is excluded by
$\re s > 0$). The construction \texttt{singlePointPacket} does not need
this analytic fact; it uses $\alpha(s)$ as a formal complex parameter
only.

\item \textbf{Singleton packet.} Build
\[
P \;:=\; \texttt{singlePointPacket}\,s\,h_{\text{Z}}
\]
of \Cref{def:single-pkt}. Its index type is $\texttt{Fin}\,1$, its
\texttt{zeroPoint} is the constant $\alpha(s)$, and its \texttt{offCritical}
flag is $\texttt{True}$. The structure is total (no missing fields) by
inspection.

\item \textbf{Index nonemptiness.} \Cref{lem:single-nonempty} gives
$h_{\text{NE}} :\texttt{Nonempty}\,P.\texttt{ZeroIndex}$ via the witness
$\langle 0\rangle:\texttt{Nonempty}\,(\texttt{Fin}\,1)$.

\item \textbf{Apply Paper 01.} By Paper 01's principal theorem
\eqref{eq:paper01},
\[
h_{\text{C}} :\texttt{Nonempty}\,
(\texttt{finiteBlaschkeModelSpace}\,P).\texttt{carrier}.
\]

\item \textbf{Choose a witness.} Apply
$\texttt{Nonempty.elim}$ (or \texttt{rcases}) to $h_{\text{C}}$ to
extract $v : (\texttt{finiteBlaschkeModelSpace}\,P).\texttt{carrier}$.
This step is constructive in the sense that any inhabitant suffices; the
specific choice is irrelevant for the statement of \texttt{Nonempty}.

\item \textbf{Apply the bridge.} Compute
\[
\texttt{bridge.transport}\,P\,h_{\text{NE}}\,v
\;:\;
\texttt{burnolBlaschkeFactor.modelSpaceCarrier}.
\]
Wrap in $\langle\,-\,\rangle$ to obtain
$\texttt{Nonempty}\,\texttt{burnolBlaschkeFactor.modelSpaceCarrier}$.
\end{enumerate}

The above is the complete chain. The Lean file
\texttt{Vol6/OffCriticalDefectKernel.lean} encodes Steps 1--7 in a single
\texttt{by} block of
\texttt{nonempty\_burnol\_modelSpaceCarrier\_of\_bridge}; the structure
projection in
\texttt{off\_critical\_zero\_defect\_kernel\_of\_bridge} then assembles
the final inhabitant.

\subsection{Lean term-mode rendering}

For reference, the Lean term has the following structure:
\begin{verbatim}
theorem nonempty_burnol_modelSpaceCarrier_of_bridge
    (bridge : OffCriticalDefectKernelBridge) :
    ExistsOffCriticalZero ->
      Nonempty burnolBlaschkeFactor.modelSpaceCarrier := by
  intro hOff
  rcases hOff with <s, hZero>
  let P : FiniteBlaschkePacket := singlePointPacket s hZero
  have hPNonempty : Nonempty P.ZeroIndex :=
    singlePointPacket_nonempty s hZero
  have hCarrier : Nonempty (finiteBlaschkeModelSpace P).carrier :=
    finite_packet_model_space_nonempty P hPNonempty
  rcases hCarrier with <v>
  exact <bridge.transport P hPNonempty v>
\end{verbatim}
The structure of the Lean tactic block is exactly the seven micro-steps,
in the same order. No tactic introduces a \texttt{sorry}, \texttt{admit},
\texttt{axiom}, or \texttt{opaque}.

\subsection{Universe analysis}

A subtle point: the bridge type
\texttt{OffCriticalDefectKernelBridge} is in \texttt{Type 1}, not in
\texttt{Prop}, because its \texttt{transport} field is genuinely a
function between arbitrary \texttt{Type}s. The conclusion
\texttt{OffCriticalZeroDefectKernel} is in \texttt{Prop}, since
\texttt{Nonempty} is a \texttt{Prop}-valued proposition. The map
\(
\texttt{OffCriticalDefectKernelBridge}\to
\texttt{OffCriticalZeroDefectKernel}
\)
is therefore a non-propositional-to-propositional function. This is
permitted in Lean 4: any proof can depend on type-level data, and a
proof of a \texttt{Prop} can be obtained by using arbitrary
\texttt{Type}-level inputs.

In the propositional view, the bridge \emph{is} the data of a
finite-packet-to-Burnol-carrier function; in the propositional view of
the conclusion, the function is consumed only to extract a
\texttt{Nonempty} witness. Universe-level coercion is automatic in Lean
because \texttt{Nonempty} is propositional and consumed only as such.

% =============================================================================
\section{Examples and Sanity Checks}\label{sec:examples}

\subsection{Why singletons suffice}

A natural objection to \Cref{def:single-pkt} is: \emph{why a singleton
packet?} Couldn't a multi-point packet be used? The answer is that
\texttt{ExistsOffCriticalZero} only supplies one zero $s$ at a time; the
existential cannot be repeated to obtain two distinct zeros without
additional information. Hence the singleton packet is the canonical
construction.

If one had access to two distinct off-critical zeros $s_{1}\neq s_{2}$,
the corresponding two-point packet
$P = \{\alpha(s_{1}), \alpha(s_{2})\}$ would yield a
\emph{two-dimensional} finite model space, with two linearly independent
Pick-basis vectors. The argument would proceed identically: Paper 01's
principal theorem produces a nonempty carrier, the bridge transports it,
and we conclude. So the proof is robust under generalisation; we use the
singleton because it matches the supplied data exactly.

\subsection{Worked example: $s = 0.6 + 14i$}

Take $s = 0.6 + 14i$ as the test point of the Haskell driver. Then
\begin{align*}
\alpha(s) &= \frac{(0.6 + 14i) - 1}{(0.6 + 14i) + 1}
            = \frac{-0.4 + 14i}{1.6 + 14i} \\
          &= \frac{(-0.4 + 14i)(1.6 - 14i)}{(1.6)^{2} + (14)^{2}}
            = \frac{(-0.64 + 196) + (5.6 + 22.4)i}{2.56 + 196} \\
          &= \frac{195.36 + 28i}{198.56}
            \approx 0.984 + 0.141i,
\end{align*}
giving $|\alpha(s)|^{2} \approx 0.968 + 0.020 = 0.988$, so
$|\alpha(s)|\approx 0.994$. The disc parameter is well inside $\Disc$ but
close to the boundary; this is consistent with $\re s = 0.6$ being close
to the critical line.

The singleton packet is
\(
P = \big(\texttt{ZeroIndex}=\texttt{Fin}\,1,\;
\texttt{zeroPoint}\,0 = 0.984+0.141i\big);
\)
its finite Blaschke product is
\(
B_{P}(z) = b_{\alpha(s)}(z)
= \frac{|\alpha(s)|}{\alpha(s)}\cdot \frac{\alpha(s)-z}{1-\overline{\alpha(s)}z}.
\)
The model space $K_{B_{P}}$ is one-dimensional, spanned by the single
Pick-basis vector
\[
k_{\alpha(s)}(z)
= \frac{1 - B_{P}(z)\overline{B_{P}(\alpha(s))}}{1 - z\overline{\alpha(s)}}.
\]
Since $B_{P}(\alpha(s)) = 0$, this simplifies to
$k_{\alpha(s)}(z) = (1 - z\overline{\alpha(s)})^{-1}$,
the standard Cauchy reproducing kernel. In Paper 01's coefficient model,
the Pick vector is the indicator function $i\mapsto[i = 0]$, which the
Haskell driver confirms is nonzero.

\subsection{Worked example: $s = 0.5 + 14i$ (on the critical line)}

For comparison, the on-line point $s = 0.5 + 14i$ has
\[
\alpha(s) = \frac{-0.5 + 14i}{1.5 + 14i}
          = \frac{(-0.5 + 14i)(1.5 - 14i)}{2.25 + 196}
          \approx 0.985 + 0.106i,
\]
so $|\alpha(s)|\approx 0.991$, very close to (but not exactly on) the
unit circle. Note that as $|\im s|\to\infty$ along the critical line,
$|\alpha(s)|\to 1$; the off-critical points sit slightly inside this
limit.

The Haskell sanity table confirms that
\texttt{isOffCriticalCandidate}$(0.5+14i) = \texttt{False}$ — \emph{not}
because the disc parameter is on the unit circle (it isn't, for finite
$s$) but because the predicate explicitly excludes
$\re s = 1/2$.

\subsection{Out-of-strip example: $s = 1.5 + 3i$}

For $s = 1.5 + 3i$ outside the critical strip, we have $\re s = 1.5$,
which violates $\re s < 1$. The Haskell predicate correctly rejects this
point. The Möbius image $\alpha(s) \approx 0.598 + 0.460i$ has
$|\alpha(s)|\approx 0.755 < 1$, well inside the disc; this is consistent
with the Cayley transform sending the entire right half-plane $\re s>0$
into $\Disc$. Such points are not zeta zeros (Mathlib's
\texttt{riemannZeta} has no zeros with $\re s > 1$), but the disc image
is still well-defined.

% =============================================================================
\section{Comparison with the Direct (Non-Singleton) Route}\label{sec:comparison}

\subsection{The direct quotient-class construction}

A naive approach to inhabiting
\texttt{burnolBlaschkeFactor.modelSpaceCarrier} would be to construct a
direct quotient class in $\HardyTwo / \BurnolFactor\HardyTwo$. This
requires:
\begin{enumerate}
\item A concrete Hardy space $\HardyTwo$ with explicit elements (Vol V's
\texttt{HardySpaceH2} is opaque);
\item A concrete Hardy submodule
$\BurnolFactor\HardyTwo$ (Vol V's \texttt{blaschkeHardyImage} is also
opaque);
\item A concrete quotient construction
$\HardyTwo / \BurnolFactor\HardyTwo$ inhabited by a specific element
not in $\BurnolFactor\HardyTwo$.
\end{enumerate}
None of these are provided by the Vol V surface. Worse, even if we
constructed them in Vol VI, the resulting carrier would be (definitionally)
some Vol VI type $X$, not Vol V's
\texttt{burnolBlaschkeFactor.modelSpaceCarrier}; we would still need a
bridge $X\to \texttt{burnolBlaschkeFactor.modelSpaceCarrier}$, returning
us to the same obstruction.

Hence the singleton packet route is preferable: it bottlenecks the
opacity exclusively at the bridge, isolating the missing surface.

\subsection{The Hardy-RKHS route}

Another conceivable route uses the Hardy-RKHS isomorphism
$\HardyTwo \simeq \mathrm{RKHS}(K_{\text{Szego}})$ where $K_{\text{Szego}}(z,w)
= (1-z\overline{w})^{-1}$ is the Szegő kernel. The model space $K_{B}$
then identifies with the RKHS subspace orthogonal to those kernels which
factor through $B$. This route is mathematically cleaner but requires Vol
V to expose the RKHS structure on $\HardyTwo$, which it does not. So this
route also reduces to the bridge.

\subsection{The completion / colimit route}

A third route lifts finite-packet model spaces to the canonical Burnol
factor by a filtered colimit
\(
\KB \;\simeq\; \mathop{\mathrm{colim}}_{P}\,K_{B_{P}}
\)
running over all finite Blaschke packets. This is the route hinted at by
RF-05 of the knowledge base. Implementing it requires both:
\begin{enumerate}
\item a category structure on finite Blaschke packets (with monomorphism
maps $P\inj P'$ when $P'$ contains $P$ as a sub-packet);
\item a colimit functor
$\mathrm{colim}: P \mapsto K_{B_{P}}$;
\item an isomorphism between this colimit and Vol V's canonical $\KB$.
\end{enumerate}
Steps (1) and (2) are routine in Mathlib; step (3) is the Vol V bridge
in disguise. So this route also bottlenecks on the same obstruction.

\subsection{Summary}

All three potential routes terminate in the same Vol V opacity. The
obstruction module captures this opacity precisely, and the conditional
principal theorem absorbs all the surrounding mathematical content into a
closed Lean theorem.

% =============================================================================
\section{The Bridge as a Mathematical Statement}\label{sec:bridge-math}

\subsection{What does it mean to inhabit the bridge?}

The bridge \texttt{OffCriticalDefectKernelBridge} asks for a function
\[
\forall P,\;\texttt{Nonempty}\,P.\texttt{ZeroIndex}\to
((\texttt{finiteBlaschkeModelSpace}\,P).\texttt{carrier}\to
\texttt{burnolBlaschkeFactor.modelSpaceCarrier}).
\]
Mathematically, this is the assertion that for any finite Blaschke
packet $P$, there exists \emph{a} map from the Pick coefficient space
$P.\texttt{ZeroIndex}\to\C$ to the canonical Burnol model space carrier.
Such a map need not be linear, isometric, or canonical; mere existence is
enough.

In particular, any one of the following Vol V upgrades would inhabit the
bridge:

\begin{enumerate}[label=(\alph*)]
\item \textbf{Concretization.} Replace
$\texttt{axiom burnolBlaschkeFactor : BurnolBlaschkeFactor}$ by a
\texttt{def} producing a concrete factor whose carrier is, say,
$\HardyTwo / \BurnolFactor\HardyTwo$ as a Lean quotient. Then the bridge
is the natural inclusion
$P.\texttt{ZeroIndex}\to\C$ \(\hookrightarrow\) $\HardyTwo / \BurnolFactor\HardyTwo$
sending a coefficient sequence to the corresponding Pick-basis sum modulo
$\BurnolFactor\HardyTwo$.

\item \textbf{Identification lemma.} Add a Vol V theorem
\(
\texttt{burnol\_modelSpaceCarrier\_eq}
:
\texttt{burnolBlaschkeFactor.modelSpaceCarrier}
=
(\HardyTwo / \BurnolFactor\HardyTwo).
\)
Then the bridge is the natural map composed with the equality cast.

\item \textbf{Constructor lemma.} Add a Vol V function
\(
\texttt{burnol\_modelSpaceCarrier\_of\_disc\_zero}
:
\forall \alpha\in\Disc,\;\BurnolFactor(\alpha)=0\to
\texttt{burnolBlaschkeFactor.modelSpaceCarrier}
\)
producing the reproducing kernel at any disc zero of $\BurnolFactor$.
Combined with a finite-packet selector, this gives the bridge.
\end{enumerate}

\subsection{Why option (a) is the intended target}

The Vol VI knowledge base RF-01 explicitly identifies (a) as the
canonical resolution: ``vol6 must define its own concrete replacement for
\texttt{burnolBlaschkeFactor} (call it
\texttt{concreteBurnolBlaschkeFactor}) that is definitionally equal or
provably equivalent.'' This is the most invasive but also the most
canonical fix; once the carrier is concrete, all three Vol VI papers
(03, 04, 05) that depend on it can be closed unconditionally.

\subsection{Why we ship the obstruction rather than the fix}

Implementing option (a) requires modifying Vol V or shipping a
significant Vol VI module that constructs the Burnol/Blaschke factor
from scratch in Lean. Both are out of scope for a single
sprint paper. Paper 03's role is to discharge \emph{Target 1} of the
no-phantom language; the carrier construction is logically prior and
should be addressed by a dedicated paper (or by a Vol V upgrade). The
obstruction module precisely localises the missing piece, so any future
worker can pick it up without re-deriving Paper 03's content.

% =============================================================================
\section{Discussion}\label{sec:discussion}

\subsection{Comparison with the Vol V opaque-axiom semantics}

Vol V's choice to declare
\texttt{burnolBlaschkeFactor} as an axiom reflects the convention that
the canonical Burnol/Blaschke factor is treated as data \emph{external}
to the Hardy/Mellin axiomatic surface used by the no-phantom route. This
keeps the core axiomatic surface lean, but it also pushes the
construction of the carrier into Vol VI. The minimal Vol V surface
upgrade that would close Paper 03 is to expose any one of:
\begin{itemize}
\item a coercion
$\texttt{HardySpaceH2}\to\texttt{burnolBlaschkeFactor.modelSpaceCarrier}$
that factors through the quotient by
$\texttt{blaschkeHardyImage}$;
\item an equivalence
$\texttt{burnolBlaschkeFactor.modelSpaceCarrier}
\simeq \texttt{HardySpaceH2}/\texttt{blaschkeHardyImage}$;
\item a constructor
$\forall \alpha\in\Disc,\;
\BurnolFactor(\alpha)=0
\to \texttt{burnolBlaschkeFactor.modelSpaceCarrier}$
producing the reproducing kernel at $\alpha$.
\end{itemize}
The third option is the most direct match for our singleton-packet
construction; the second is the most categorically natural; the first is
the most immediate Hardy-space-theoretic statement.

\subsection{Why this paper is not RH-equivalent}

Paper 03 \emph{conditionally} discharges Target 1 of the no-phantom
language. Combined with Target 2 (Vol VI Paper 06's discharge of
\texttt{YonedaBlaschkeDetectorCalculus}) and the Vol V master theorem
\eqref{eq:master}, an unconditional Paper 03 closes
\texttt{RH\_classical}. We must therefore explain why our argument is
not circular and not RH-equivalent.

The hypothesis \texttt{ExistsOffCriticalZero} of
\Cref{thm:cond-principal} is the \emph{negation} of RH; the conclusion is a
nonemptiness statement about the canonical Burnol model space. The
implication, properly read, is: \emph{any} off-critical zero (whether or
not RH holds) creates a nonzero model-space vector. The Vol V master
theorem then turns that contradiction into a proof of RH only via
combination with the no-phantom predicate
$\texttt{NoPhantomDefect}$, which is supplied by Target 2 (Paper 06).
Neither Paper 03 nor Paper 06 individually proves RH; their conjunction
under the Vol V master rule does. This is the new-language pattern: the
mathematical content has been factored into two strictly weaker pieces,
each provable by direct construction.

\subsection{Limitations}

\begin{enumerate}
\item \textbf{Bridge inhabitation.} \Cref{thm:cond-principal} is
conditional on the bridge. The remaining mathematical work for an
unconditional Vol VI closure is the inhabitation of
\texttt{OffCriticalDefectKernelBridge}, equivalent to discharging RF-01.
\item \textbf{Paper 01 carrier choice.} Paper 01's choice of carrier
$P.\texttt{ZeroIndex}\to\C$ is the simplest valid finite-dimensional
carrier; it reproduces the Pick basis up to a noncanonical isomorphism.
Other carriers (e.g.\ a Hardy-space subspace generated by
$\{k_{\alpha_{i}}\}$) would also work and might make the bridge
inhabitation easier; the present obstruction is parametric in the carrier
and remains valid for any Paper 01 implementation.
\item \textbf{Numerical scope.} The Haskell program is a single-point
sanity check, not a verification of any specific zeta zero.
\end{enumerate}

% =============================================================================
\section{Conclusion}\label{sec:conclusion}

We have discharged Step~(I) and Step~(II) of the Paper 03 pipeline as
closed Lean theorems
\begin{itemize}
\item $\texttt{singlePointPacket\_nonempty}$
   (\Cref{lem:single-nonempty}),
\item $\texttt{singlePointModelSpace\_nonempty}$
   (\Cref{thm:single-nonempty}),
\item $\texttt{finitePacketWitness\_of\_offcritical}$
   (\Cref{cor:fin-witness}),
\item $\texttt{exists\_finitePacket\_witness\_of\_offCritical}$
   (\Cref{cor:exists-witness}),
\end{itemize}
and we have shipped Step~(III) as the conditional principal theorem
\(
\texttt{off\_critical\_zero\_defect\_kernel\_of\_bridge}
:
\texttt{OffCriticalDefectKernelBridge}\to
\texttt{OffCriticalZeroDefectKernel}
\)
together with the typed obstruction module
\texttt{Vol6.Obstruction.OffCriticalDefectKernelObstruction}. The
unconditional principal theorem
\texttt{off\_critical\_zero\_defect\_kernel} is \emph{not} declared in
the present module; doing so without a Vol V bridge would violate the
Vol VI hard rules.

The Lean development compiles as
\(
\texttt{lake build Vol6.OffCriticalDefectKernel}
\)
with no \texttt{sorry}, \texttt{admit}, new \texttt{axiom}, or new
\texttt{opaque}. The accompanying Haskell program runs successfully and
mirrors Step~(I) and Step~(II) numerically.

\paragraph{Future work.} The next concrete step in the Vol VI pipeline
is the production of an inhabitant of
\texttt{OffCriticalDefectKernelBridge}. This requires either a Vol V
upgrade exposing one of the three bridge surfaces listed in
\Cref{sec:transport-obstruction}, or a Vol VI replacement of
\texttt{burnolBlaschkeFactor} by a concrete (non-axiomatic) Burnol
factor. The latter is anticipated by RF-01 and would be a substantive
mathematical contribution: a Lean-formalised construction of the
Burnol/Blaschke factor in terms of fractional-part atoms and
finite-rank Hardy quotients.

% =============================================================================
\section{Appendix: Full Lean Module Listing}\label{sec:lean-listing}

For completeness we reproduce the full content of
\texttt{lean/Vol6/OffCriticalDefectKernel.lean} (the principal module) and
\texttt{lean/Vol6/Obstruction/OffCriticalDefectKernelObstruction.lean}
(the obstruction module). These are the canonical Lean artefacts of
Paper 03 and should be regarded as part of the paper.

\subsection{Principal module}

Located at \texttt{lean/Vol6/OffCriticalDefectKernel.lean}. The module
namespace is \texttt{Vol6.OffCriticalDefectKernelNS} and the
\texttt{noncomputable section} mode is used (because elements are
extracted via \texttt{Classical.choice} from \texttt{Nonempty}).

Key declarations:
\begin{itemize}
\item \texttt{abbrev alpha (s : Complex) : Complex := mobiusParam s} — disc parameter map.
\item \texttt{def singlePointPacket (s : Complex) (\_h : OffCriticalZero s) : FiniteBlaschkePacket}
   — singleton packet construction.
\item \texttt{theorem singlePointPacket\_nonempty} — closed.
\item \texttt{def singlePointModelSpace} — singleton model space.
\item \texttt{theorem singlePointModelSpace\_nonempty} — closed via
   Paper 01.
\item \texttt{theorem nonempty\_burnol\_modelSpaceCarrier\_of\_bridge}
   — Step~(III) under the bridge.
\item \texttt{theorem off\_critical\_zero\_defect\_kernel\_of\_bridge}
   — conditional principal theorem.
\item \texttt{theorem finitePacketWitness\_of\_offcritical} — audit
   corollary.
\item \texttt{theorem exists\_finitePacket\_witness\_of\_offCritical}
   — audit corollary.
\item \texttt{def off\_critical\_defect\_kernel\_status : String}
   — audit marker.
\end{itemize}

\subsection{Obstruction module}

Located at \texttt{lean/Vol6/Obstruction/OffCriticalDefectKernelObstruction.lean}.
The module namespace is \texttt{Vol6.Obstruction}.

Key declarations:
\begin{itemize}
\item \texttt{noncomputable def mobiusParam (s : Complex) : Complex := (s - 1)/(s + 1)}
\item \texttt{theorem mobiusParam\_one : mobiusParam 1 = 0}
\item \texttt{def FiniteBlaschkePacketCarrierTransport : Type 1}
\item \texttt{structure OffCriticalDefectKernelBridge : Type 1}
\item \texttt{def SingletonBlaschkePacketCarrierTransport : Type 1}
\item \texttt{def singleton\_transport\_of\_full\_transport}
\item \texttt{def off\_critical\_defect\_kernel\_obstruction\_status : String}
\end{itemize}

\subsection{Build verification}

The Lean development is verified by
\begin{verbatim}
cd lean && PATH="$HOME/.elan/bin:$PATH" \
    lake build Vol6.OffCriticalDefectKernel
\end{verbatim}
which terminates with
\begin{verbatim}
Build completed successfully (2925 jobs).
\end{verbatim}
A grep for \texttt{sorry|admit|axiom |opaque } in
\texttt{Vol6/OffCriticalDefectKernel.lean} and
\texttt{Vol6/Obstruction/OffCriticalDefectKernelObstruction.lean}
yields only matches in comments (the words appear in audit prose
explaining what is forbidden).

% =============================================================================
\section{Appendix: Full Haskell Driver Listing}\label{sec:haskell-listing}

The Haskell program lives at
\texttt{haskell/paper-03-off-critical-defect-kernel/Main.hs}. Key
components:

\begin{itemize}
\item \texttt{mobiusParam :: Complex Double -> Complex Double} —
   numerical Möbius parameter.
\item \texttt{isOffCriticalCandidate :: Complex Double -> Bool}
   — strip and off-line predicate.
\item \texttt{data SinglePointPacket} — single-zero packet.
\item \texttt{type SinglePointCarrier = () -> Complex Double} — singleton
   carrier (mirroring \texttt{Fin 1 -> Complex}).
\item \texttt{singlePointKernelVector :: SinglePointCarrier} — the
   constant-1 function.
\item \texttt{kernelVectorNonzero :: Bool} — the nondegeneracy decision.
\item \texttt{runPipeline :: Complex Double -> IO Bool} — full driver.
\item \texttt{sanityTable :: IO ()} — multi-point smoke test.
\end{itemize}

The program runs successfully with
\begin{verbatim}
runghc haskell/paper-03-off-critical-defect-kernel/Main.hs
\end{verbatim}
producing output:
\begin{verbatim}
Vol6 paper 03 — Off-Critical Defect Kernel (Haskell)
Input s = 0.6 :+ 14.0
  Re s    = 0.6
  Im s    = 14.0
alpha(s) = (s-1)/(s+1) = 0.984 :+ 0.141
  |alpha(s)|  = 0.994
Off-critical candidate? True
Singleton kernel vector v : () -> Complex
  v()      = 1.0 :+ 0.0
  |v|      = 1.0
  nonzero? = True
Assertions:
  [A1] off-critical strip + off line .... True
  [A2] 0 < |alpha(s)| < 1 ............... True
  [A3] zeroPoint() == alpha(s) .......... True
  [A4] kernel vector norm > 0 ........... True
  [A5] kernelVectorNonzero flag ......... True
RESULT: PASS
\end{verbatim}

% =============================================================================
\section{Appendix: Vol V Surface Audit}\label{sec:vol5-audit}

We document the precise Vol V audit performed to certify that no
introduction rule for \texttt{burnolBlaschkeFactor.modelSpaceCarrier}
exists in the no-phantom-language route.

\subsection{Files searched}

\begin{itemize}
\item \texttt{Vol5/YonedaNymanTrackA/BurnolCore.lean}
\item \texttt{Vol5/YonedaNymanTrackA/BurnolFactorization.lean} (excluded
from the no-phantom route per knowledge base RF-06; surveyed for
completeness only)
\item \texttt{Vol5/YonedaNymanTrackA/BlaschkeDefectObject.lean}
\item \texttt{Vol5/YonedaNymanTrackA/HardyModelSpace.lean}
\item \texttt{Vol5/YonedaNymanTrackA/NoPhantomConcreteAtoms.lean}
\item \texttt{Vol5/YonedaNymanTrackA/NoPhantomBlaschkeDefect.lean}
\item \texttt{Vol5/YonedaNymanTrackA/CondensedHilbertDefect.lean}
\item \texttt{Vol5/YonedaNymanTrackA/RKHSDetectorSemantics.lean}
\item \texttt{Vol5/YonedaNymanTrackA/FiniteRankRKHSDetector.lean}
\item \texttt{Vol5/YonedaNymanTrackA/OffCriticalZeroDefect.lean}
\item \texttt{Vol5/YonedaNymanTrackA/NoPhantomLanguage.lean}
\item \texttt{Vol5/YonedaNymanTrackA/YonedaCore.lean}
\item \texttt{Vol5/YonedaNymanTrackA/Basic.lean}
\end{itemize}

\subsection{Search method}

For each file, we ran (essentially)
\begin{verbatim}
grep -n "modelSpaceCarrier" file.lean
\end{verbatim}
and inspected each match for context. The matches in
\texttt{BurnolCore.lean}, \texttt{BlaschkeDefectObject.lean},
\texttt{NoPhantomConcreteAtoms.lean}, and \texttt{NoPhantomLanguage.lean}
are uses, not introductions: they reference the carrier via the axiom
\texttt{burnolBlaschkeFactor.modelSpaceCarrier} and either project from
it or assert membership/emptiness predicates. None constructs an
inhabitant.

\subsection{Conclusion}

No Vol V theorem, function, or instance has any of the conclusion
shapes (a) identification, (b) injection, or (c) constructor for
\texttt{burnolBlaschkeFactor.modelSpaceCarrier}, for any concrete type.
This certifies that the obstruction module is necessary; without a Vol V
upgrade, the unconditional principal theorem cannot be closed by purely
axiom-free Vol VI means.

\subsection{Implications for the Vol VI sprint}

This audit shows that the Vol VI Paper 03 obstruction is not unique:
Papers 04, 05, and 06 likely face analogous obstructions for the
\texttt{DefectDetectedByRationalDilationRepresentable} (RF-02) and
\texttt{BlaschkeDefectDualizable/Polarized/Regularized} (RF-03)
opaque atoms. Each of these papers will need to ship a parallel
obstruction module, and the synthesis paper (Paper 07) will need to
collect and propagate them. The pattern is uniform: when a Vol V
opaque/axiom can be inhabited only externally, Vol VI factors out a
typed bridge and delivers conditional theorems. The conditional
theorems are then assembled by the orchestrator pipeline; the
obstruction propagation tracks exactly what Vol V upgrades remain
required for full closure of \texttt{RH\_classical\_new\_language}.

% =============================================================================
\section*{Acknowledgements}

We thank the orchestrator-pipeline of the Vol VI sprint for parallel
delivery of Vol VI Paper 01 (\texttt{Vol6.FiniteBlaschkePacket}), without
which this paper's Step~(II) would have to be re-derived in situ. The
Vol V no-phantom-language route, on which all of Vol VI rests, is due to
the Vol V development team.

% =============================================================================
\begin{thebibliography}{99}

\bibitem{burnol}
J.-F. Burnol, \emph{Sur certains espaces de Hilbert de fonctions
enti\`eres, li\'es \`a la transformation de Fourier et aux fonctions L
de Dirichlet et de Riemann}, C. R. Acad. Sci. Paris S\'er. I Math.,
333 (2001), 273--278.

\bibitem{nyman}
B. Nyman, \emph{On the One-Dimensional Translation Group and
Semi-Group in Certain Function Spaces}, Ph.D. thesis, University of
Uppsala, 1950.

\bibitem{beurling}
A. Beurling, \emph{A closure problem related to the Riemann
zeta-function}, Proc. Nat. Acad. Sci. U.S.A., 41 (1955), 312--314.

\bibitem{garcia-mashreghi}
S. R. Garcia, J. Mashreghi, W. T. Ross, \emph{Introduction to Model
Spaces and Their Operators}, Cambridge Studies in Advanced
Mathematics, Vol. 148, Cambridge University Press, 2016.

\bibitem{hoffman}
K. Hoffman, \emph{Banach Spaces of Analytic Functions}, Prentice-Hall,
Englewood Cliffs, NJ, 1962.

\bibitem{nikolski}
N. K. Nikolski, \emph{Operators, Functions, and Systems: An Easy
Reading}, Vol. 1: Hardy, Hankel, and Toeplitz, Mathematical Surveys
and Monographs, Vol. 92, American Mathematical Society, 2002.

\bibitem{pick}
G. Pick, \emph{\"Uber die Beschr\"ankungen analytischer Funktionen},
Math. Ann., 77 (1916), 7--23.

\bibitem{rudin}
W. Rudin, \emph{Real and Complex Analysis}, 3rd ed., McGraw-Hill, New
York, 1987.

\bibitem{titchmarsh}
E. C. Titchmarsh, \emph{The Theory of the Riemann Zeta-Function}, 2nd
ed., Oxford University Press, 1986.

\bibitem{vol5}
M. Long, \emph{The No-Phantom Language for Classical RH}, Vol V of the
HoTT--Riemann programme, Magneton Research, 2025.

\bibitem{vol6-knowledge-base}
M. Long, \emph{Vol VI Knowledge Base}, Magneton Research,
\texttt{hott-riemann-hypothesis-vol6/.knowledge-base.md}, 2026.

\bibitem{vol6-paper01}
M. Long, \emph{Concrete Finite Blaschke Packet Model Spaces} (Vol VI
Paper 01), Magneton Research, 2026.

\bibitem{vol6-paper02}
M. Long, \emph{Detector Completeness for Finite Blaschke Packets via
Gram-Matrix Nondegeneracy} (Vol VI Paper 02), Magneton Research, 2026.

\bibitem{vol6-paper06}
M. Long, \emph{Assembly of the Yoneda/Blaschke Detector Calculus} (Vol
VI Paper 06), Magneton Research, 2026.

\bibitem{lean4}
L. de Moura, S. Ullrich, \emph{The Lean 4 Theorem Prover and
Programming Language}, in: A. Platzer, G. Sutcliffe (eds.), Automated
Deduction --- CADE 28, Lecture Notes in Computer Science, vol. 12699,
Springer, Cham, 2021, pp. 625--635.

\bibitem{mathlib}
The Mathlib Community, \emph{The Lean Mathematical Library}, in:
Proceedings of the 9th ACM SIGPLAN International Conference on
Certified Programs and Proofs (CPP 2020), New Orleans, 2020,
pp. 367--381.

\bibitem{voevodsky}
V. Voevodsky and the Univalent Foundations Program, \emph{Homotopy
Type Theory: Univalent Foundations of Mathematics}, Institute for
Advanced Study, Princeton, 2013.

\end{thebibliography}

\end{document}
